\begin{abstract}
Model merging has emerged as a powerful, training-free paradigm for combining task-specific experts into a single unified model. Recently, dynamic model merging has gained prominence by predicting blending coefficients on-the-fly for each input sample. However, state-of-the-art dynamic methods have introduced highly complex, over-engineered architectures---such as quantum wavefunction superposition and high-dimensional phase projectors---arguing that classical linear routing is structurally limited and prone to catastrophic representation collapse on challenging tasks like SVHN. In this paper, we deconstruct this assumption through the lens of Occam's razor. We demonstrate that the failure of classical linear routing is not a structural limitation, and that a properly configured classical unregularized Linear Router is already highly robust in standard homogeneous settings (achieving $91.53\% \pm 0.41\%$ mean accuracy across 5 seeds). To further address out-of-distribution shifts and heterogeneous deployment scenarios, we propose \textbf{Robust Linear Routing (RLR)}, which regularizes a simple 768-parameter classical gating layer using standard $L_2$ weight decay and Softmax Temperature scaling. Rather than being a mandatory fix for homogeneous collapse, RLR acts as a specialized stabilizer that trades off a negligible margin of peak homogeneous performance to secure superior resilience under mixed-task heterogeneous test streams, preventing batch-level coefficient collapse. Our findings show that simple classical regularization is all we need to stabilize and protect dynamic model merging under out-of-distribution and heterogeneous settings, challenging the necessity of needlessly complex parameter fusion frameworks.
\end{abstract}
